% MODEL-ID: SP-P0002-TTNPU-Q36-CP-v1
% UPSTREAM-COMMIT: 7864d5dc17930667d663bbadd1ce2bc722de2753
% CP-DEGREE: 4
% GLOBAL-SEQ-LEN: 32768
% LOCAL-SEQ-LEN: 8192
% FULL-Q-HEADS: 24
% FULL-KV-HEADS: 4
% GDN-QK-HEADS: 16
% GDN-V-HEADS: 48

\documentclass[11pt,a4paper]{article}
\usepackage[T1]{fontenc}
\usepackage{lmodern}
\usepackage{amsmath,amssymb,mathtools,booktabs,geometry,hyperref,microtype}
\geometry{margin=27mm}
\hypersetup{colorlinks=true,linkcolor=black,urlcolor=blue,citecolor=black}
\newcommand{\CP}{\operatorname{CP}}
\newcommand{\Tmap}{\mathcal T}
\newcommand{\Kernel}{\mathcal K}
\title{Conditional Correctness of Qwen3.6 Context Parallel in TorchTitan-NPU \texttt{override-refactor}}
\author{Model ID: SP-P0002-TTNPU-Q36-CP-v1}
\date{2026-08-24}

\begin{document}
\maketitle

\begin{abstract}
We formalize the Context Parallel (CP) transport used by the Qwen3.5/3.6
long-text path in the \texttt{override-refactor} branch of
\texttt{cann/torchtitan-npu}, frozen at commit
\texttt{7864d5dc17930667d663bbadd1ce2bc722de2753}. The reviewed configuration
shards a sequence of length $32768$ over CP degree $4$, leaving $8192$ tokens
on each sequence rank. It applies CP both to variable-length TND full attention
and to Gated DeltaNet (GDN), which includes causal convolution and recurrent
state updates.

The proof separates transport, local operator semantics, and the code-to-model
bridge. Lean 4 verifies that changing from sequence sharding to head-block
sharding is a rank-axis transpose, that the inverse transform restores the
original tensor, and that packing several projected tensors before the
collective is componentwise transparent. Consequently, if the real collective
implements this lossless transpose, if the local fused kernel implements the
same head-block-separable function as the non-CP baseline, and if both paths
receive identical global sequence metadata, then the CP output is pointwise
equal to the dense output. The theorem covers full attention and GDN, but does
not certify DTensor, HCCL, CANN kernels, metadata construction, floating-point
rounding, or distributed backward.
\end{abstract}

\section{External Evidence and Reviewed Scope}

The reviewed source paths are:
\begin{verbatim}
torchtitan_npu/models/qwen3_5/config_registry.py
torchtitan_npu/override/qwen3_5/parallelize.py
torchtitan_npu/override/qwen3_5/varlen_attention.py
torchtitan_npu/override/qwen3_5/gated_delta.py
\end{verbatim}
The Python package remains named \texttt{qwen3\_5}, while the override
description explicitly refers to ``Qwen3.5/3.6 Context Parallel.'' We retain
the exact paths and use the user-facing name Qwen3.6.

The long-text configuration fixes
\[
L=32768,\qquad P=4,\qquad L_{\mathrm{local}}=8192.
\]
It loads both the variable-length attention CP override and the GDN CP
override.

\subsection{Transport Skeleton}

\texttt{sequence\_to\_head\_shard} declares a local tensor as sharded along
the sequence dimension and redistributes it to head sharding.
\texttt{head\_to\_sequence\_shard} performs the reverse redistribution.
\texttt{exchange\_sequence\_heads} packs head chunks from several projections
before communication. Packing may reduce collective calls, but it must not
change ownership or order of elements.

\subsection{Full Attention}

The full-attention path performs optional GQA K/V handling, exchanges Q/K/V
from sequence sharding to head-block sharding, invokes TND fused attention on
the complete sequence, and exchanges the output back. The reviewed head counts
are
\[
H_Q=24,\qquad H_{KV}=4.
\]
With $P=4$, each rank receives
\[
H_Q^{\mathrm{local}}=6,\qquad H_{KV}^{\mathrm{local}}=1.
\]
Thus one local KV head aligns with six local Q heads.

\subsection{Gated DeltaNet}

The GDN path exchanges Q, K, V, decay, and beta. It then executes
segment-aware causal convolution and recurrent updates over the complete
sequence, before returning the output to sequence sharding. Its head counts
are
\[
H_{QK}=16,\qquad H_V=48,
\]
which become 4 and 12 heads per CP rank. The formal model treats the entire
contiguous head chunk owned by one rank as an abstract payload, rather than
pretending that all projections have the same scalar-head shape.

\section{Mathematical Model and Assumptions}

Represent a global sequence position by
\[
(r_s,i),\qquad r_s\in\{0,\ldots,P-1\},\quad
i\in\{0,\ldots,L_{\mathrm{local}}-1\},
\]
and a global head block by
\[
r_h\in\{0,\ldots,P-1\}.
\]
In sequence-sharded layout, write an element as $X[r_s,i,r_h]$. Define the
sequence-to-head transform by
\[
(\Tmap X)[r_h,r_s,i]=X[r_s,i,r_h], \tag{1}
\]
and its inverse by
\[
(\Tmap^{-1}Y)[r_s,i,r_h]=Y[r_h,r_s,i]. \tag{2}
\]

\paragraph{Assumption A: lossless transport.}
The actual collective implements (1) without loss, duplication, local token
reordering, or incorrect chunk boundaries.

\paragraph{Assumption B: local-kernel semantic agreement.}
For each head block $r_h$, the CP fused kernel and the dense baseline implement
the same complete-sequence function
\[
\Kernel_m(Q_{r_h},K_{r_h},V_{r_h}),
\]
where $m$ contains causal boundaries, cumulative sequence lengths, reset
markers, and scaling information. The operator may be full attention or an
entire GDN head-block pipeline.

\paragraph{Assumption C: identical metadata.}
The CP path reconstructs exactly the same global segmentation as the baseline.
It neither joins distinct packed samples into one causal segment nor omits a
GDN reset.

\section{General Equivalence Theorem}

Define the dense head-separable result by
\[
Y_{\mathrm{dense}}[r_s,i,r_h]
=
\Kernel_m
\left(
Q[\cdot,\cdot,r_h],
K[\cdot,\cdot,r_h],
V[\cdot,\cdot,r_h]
\right)[r_s,i]. \tag{3}
\]
Define the CP result by
\[
Y_{\mathrm{cp}}
=
\Tmap^{-1}
\left(
r_h\mapsto
\Kernel_m
\left(
(\Tmap Q)[r_h,\cdot,\cdot],
(\Tmap K)[r_h,\cdot,\cdot],
(\Tmap V)[r_h,\cdot,\cdot]
\right)
\right). \tag{4}
\]

\subsection{Round Trip}

\textbf{Theorem 1.}
For every tensor $X$,
\[
\Tmap^{-1}(\Tmap X)=X,\qquad
\Tmap(\Tmap^{-1}X)=X. \tag{5}
\]

\textbf{Proof.}
Fix $(r_s,i,r_h)$. By definition,
\[
(\Tmap^{-1}\Tmap X)[r_s,i,r_h]
=(\Tmap X)[r_h,r_s,i]
=X[r_s,i,r_h].
\]
The opposite direction is identical. Lean applies function extensionality to
the three indices; the remaining proof reduces to \texttt{rfl}. \qed

\subsection{CP Forward Equals Dense Forward}

\textbf{Theorem 2.}
For arbitrary Q, K, V, metadata, and any head-block-separable kernel,
\[
\boxed{Y_{\mathrm{cp}}=Y_{\mathrm{dense}}}. \tag{6}
\]

\textbf{Proof.}
Fix $(r_s,i,r_h)$. Expanding (4) and the inverse transform gives
\[
Y_{\mathrm{cp}}[r_s,i,r_h]
=
\Kernel_m
\left(
(\Tmap Q)[r_h,\cdot,\cdot],
(\Tmap K)[r_h,\cdot,\cdot],
(\Tmap V)[r_h,\cdot,\cdot]
\right)[r_s,i].
\]
Equation (1) yields
\[
(\Tmap Q)[r_h,r'_s,i']=Q[r'_s,i',r_h],
\]
and similarly for K and V. The right-hand side is exactly (3). Since the
output index was arbitrary, function extensionality gives equality of the
whole tensors. In Lean, the indexed equality reduces to \texttt{rfl}. \qed

\subsection{Packed Communication}

Let
\[
P(Q,K,V)[r_s,i,r_h]
=
(Q[r_s,i,r_h],K[r_s,i,r_h],V[r_s,i,r_h]).
\]
Then
\[
\Tmap(P(Q,K,V))
=
P(\Tmap Q,\Tmap K,\Tmap V). \tag{7}
\]
Therefore packing projections before the collective is semantically
transparent whenever the physical chunk and split boundaries match the
payload partition.

\subsection{Downstream Observations}

For any pure downstream observation $\Phi$,
\[
\Phi(Y_{\mathrm{cp}})=\Phi(Y_{\mathrm{dense}}). \tag{8}
\]
Lean obtains this by congruence from (6). If the two sides are interpreted as
one differentiable real function, their mathematical derivatives agree where
they exist. A concrete autograd engine, BF16 rounding, and distributed
gradient collectives are outside the formal model.

\section{Full-Attention Instance}

At one head-block rank, the Q payload may contain six consecutive Q heads,
while the K and V payloads each contain one KV head. The local operator
contains the GQA mapping. Alignment follows from
\[
24=4\times6,\qquad
4=4\times1,\qquad
6=1\times6.
\]
Lean checks these closed integer equalities. Causal masking, TND layout,
variable-length boundaries, and scale remain in metadata or in the local
kernel. The transport theorem says CP itself does not alter these
mathematical inputs when the complete token order is preserved.

\section{GDN Instance}

Define one compound GDN payload
\[
Z=(Q,K,V,a,\beta,\theta_{\mathrm{conv}},A_{\log},b_{\Delta t}),
\]
and let $\mathcal G_m(Z)$ perform segment-aware causal convolution, decay-gate
construction, reset-aware Gated Delta recurrence, and output generation. The
unary instance of the theorem is
\[
\Tmap^{-1}\bigl(\mathcal G_m(\Tmap Z)\bigr)
=
\mathcal G_m(Z). \tag{9}
\]
The source-level bridge must still show that \texttt{shard\_local\_heads}
aligns convolution weights, \texttt{A\_log}, and \texttt{dt\_bias} with the
exchanged activation block, and that cumulative sequence lengths induce the
same resets as the baseline.

\section{Formal Verification}

\texttt{Proof.lean} imports only Lean Core \texttt{Init}. It verifies:
\begin{itemize}
  \item $32768=4\times8192$;
  \item the 24/4 full-attention split into 6/1 heads per rank;
  \item the 16/48 GDN split into 4/12 heads per rank;
  \item both rank-axis transforms are mutual inverses;
  \item packed and componentwise transport are equal;
  \item CP/dense equality for arbitrary metadata and arbitrary local kernels;
  \item preservation of arbitrary pure downstream observations.
\end{itemize}

The intended audit commands are:
\begin{verbatim}
lake build
lake env leanchecker --fresh Problems.TorchTitanNPUQwen36CP.Proof
lake env lean Problems/TorchTitanNPUQwen36CP/Proof.lean
\end{verbatim}
The file uses no \texttt{sorry}, \texttt{admit}, custom axiom, or
\texttt{native\_decide}; \texttt{\#print axioms} exposes dependencies.

\section{Interpretation Boundary}

The theorem does not prove that DTensor/HCCL can never reorder data in a
specific software stack. It does not prove every dynamic
\texttt{torch.chunk}/concatenate/split boundary, every packed or padded
metadata case, bitwise BF16 equivalence of CANN attention, or bitwise
equivalence of the Triton-Ascend GDN forward, backward, and checkpoint replay.
It also does not cover independent defects in FSDP, optimizer updates, or
checkpoint serialization.

Engineering validation should therefore include deterministic CP1/CP4
comparisons from the same checkpoint and global batch, intermediate Q/K/V and
GDN-state comparisons, cross-rank packed-sequence boundary cases, randomized
transport inversion tests, 32K/64K loss and gradient comparisons, and
save/reload continuation tests.

\section{Conclusion}

The reviewed Qwen3.6 CP skeleton is
\[
\text{sequence shard}
\xrightarrow{\Tmap}
\text{head shard with full sequence}
\xrightarrow{\Kernel_m}
\text{head-sharded output}
\xrightarrow{\Tmap^{-1}}
\text{sequence-sharded output}.
\]
Lean verifies
\[
\boxed{\Tmap^{-1}\circ\Kernel_m\circ\Tmap=\Kernel_m}
\]
pointwise for every head-block-separable $\Kernel_m$ with identical metadata.
The defensible conclusion is that the CP communication skeleton is
mathematically correct under the stated transport, kernel, and metadata
premises. Remaining risk lies in the code-to-model correspondence, especially
DTensor chunk boundaries, variable-length metadata, fused-kernel semantics,
and the backward/training path.

\end{document}
